\section{Introduction}
\label{sec:intro}

A seasoned trader's edge is not programmed; it is earned through thousands of profit-and-loss cycles that gradually reshape judgment about which geometric configurations on a chart warrant action, which factor combinations confirm a thesis, and how aggressively to size a position under a given volatility regime. Each realized outcome feeds back into the trader's internal model, updating both parametric intuitions and explicit playbook entries. This experiential loop is precisely what current financial multi-agent systems lack.

Recent LLM-based financial agents have made impressive strides in role decomposition and collaborative reasoning~\cite{tradingagents2025,fincon2024,finagent2024}. Yet their self-improvement mechanisms remain confined to \emph{verbal reinforcement}: reflections on past trades are summarized into natural-language beliefs and prepended to future prompts. Risk-adjusted returns, conditional value-at-risk, and maximum drawdown, the quantitative signals that professional traders internalize over years, never reach the model's parameters. A live benchmark covering twelve months of real trading reveals that even frontier closed-source models incur significant net losses~\cite{deepfund2025}; separately, reinforcement-finetuned LLM agents achieve higher directional accuracy yet lower cumulative returns~\cite{losingwinner2025}, exposing a reward-hacking failure mode when the optimization target misaligns with portfolio-level objectives.

These observations point to a fundamental gap we term the \emph{evolution deadlock}: without a mechanism that propagates trajectory-level rewards back to both model parameters and long-term memory, financial agents cannot adapt to regime shifts and will inevitably suffer alpha decay. Verbal reinforcement touches neither the parametric nor the memory consolidation pathway; it is the equivalent of a trader who writes journal entries but never updates muscle memory.

Meanwhile, the broader LLM community has witnessed rapid progress in on-policy self-distillation as a post-training paradigm. Self-Distilled Reasoner~\cite{opsd2026} demonstrates that a single model can serve as both student and teacher through differential conditioning, where the teacher observes privileged information (ground-truth solutions) that the student must learn to approximate. Token-level distribution matching via Jensen-Shannon divergence yields substantial gains in mathematical reasoning without requiring an external teacher or multiple rollouts per problem. Concurrent works on iterative self-policy distillation~\cite{crisp2026} and trajectory-level tree search~\cite{seea2025} confirm the generality of this paradigm.

We ask: can on-policy self-distillation be the mechanism that breaks the evolution deadlock in financial multi-agent systems? The transfer is non-trivial. Financial trajectories are sparse and high-variance compared to math problems; credit must be assigned across multiple specialist agents rather than a single reasoner; and the reward signal (risk-adjusted PnL) is noisy, delayed, and regime-dependent. Addressing these challenges requires domain-specific adaptations: Shapley-based credit sampling to disentangle agent contributions, per-token KL clipping to prevent stylistic tokens from dominating gradients, and a complementary non-parametric memory track that consolidates stable winning patterns into retrievable belief entries.

We present \textsc{FinOPD}, a financial multi-agent framework whose core contribution is an experience-driven dual-track self-evolution mechanism. On the parametric side, the same LoRA-adapted model simultaneously defines a student policy (conditioned on current observations) and a teacher policy (additionally conditioned on hindsight risk-adjusted returns as privileged information). Token-level JSD matching drives the student toward the teacher's informed distribution, with gradients weighted by Shapley credit across five specialist agents. On the non-parametric side, winning geometry-factor-action triples are extracted from high-reward trajectories and indexed by visual-geometric embeddings for retrieval-augmented inference. Together, these two tracks convert trading experience into sustained capability improvement along both the parametric and memory axes.

Our contributions are as follows. (1) We introduce on-policy self-distillation to financial multi-agent systems, using hindsight risk-adjusted returns as privileged information and Shapley credit assignment to handle multi-agent trajectories. (2) We design a complementary non-parametric belief consolidation mechanism that indexes winning decision triples by geometric embeddings for retrieval-augmented inference. (3) We demonstrate that the resulting evolution curve exhibits observable monotonic improvement across iterations under strict post-cutoff evaluation, providing the first empirical evidence that policy-level self-evolution in financial agents is both achievable and measurable. Across Dow-30 and CSI-300 benchmarks, \textsc{FinOPD} achieves consistent improvements over twelve baselines in Sharpe ratio, maximum drawdown, and cross-regime robustness.
